\chapter{Combinatoria}
En esta sección exploraremos los conceptos
fundamentales de la combinnatoria.

Primero hablaremos de conceptos más
generales.

\section{Propiedades sumativas}

\section{Coloreo}

\section{Casos}

\section{Insight}

\section{Tableros}

\section{}

\section{Problemas}
\begin{problem}[El Salvador 2016]
Se tienen $15$ fichas circulares,
todas del mismo radio en una formación de triángulo
equilátero como se muestra en la figura.
Cada una de las fichas se pinta de un solo color:
rojo o azul.

Demuestre que, sin importar cómo se pinten las fichas,
siempre existen tres de ellas del mismo color
cuyos centros son los vértices de un triángulo
equilátero.

\begin{hints}
	\begin{hint}
		Considera qué pasa con el triángulo central.
	\end{hint}

	\begin{hint}
		Sin pérdida de generalidad dos fichas del triángulo
		central son rojas y la otra es azul.
	\end{hint}

	\begin{hint}
		Descartar casos.
	\end{hint}
\end{hints}
\end{problem}

\centering
\begin{asy}
	size(8cm);
	for(real i = 0; i < 5; ++i){
			for(real j = -0.5 * i; j <= 0.5*i; ++j){
					draw(circle((j, -i * sqrt(3) * 0.5), 0.5));
				}
		}
\end{asy}
\par
\raggedright

\begin{problem}
En un torneo de fútbol participan $30$ equipos, cada equipo
juega contra cada otro exactamente dos veces. Calcule el número
de encuentros del torneo.
\end{problem}

\begin{problem}
Considere que cada punto en el plano euclídeo está coloreado de rojo o verde.
Demuestre que existen dos puntos a distancia exactamente $1$ que son del mismo color.
\end{problem}

